\section{Allowing hypotheses on supplied lines}
\label{sec:boundaries}

Allowing a hypothesis to lie on a queried line adds a third outcome.
To compare with the binary algorithms, encode each line answer using
at most two bits: the first separates zero from nonzero, and the
second separates negative from positive on the nonzero branch. Thus
each ternary query expands to at most two levels in a binary comparison
tree, which accounts for the factor of two in the guarantee.

Algorithm~\ref{alg:boundary} implements the reduction below.
In this section a supplied line returns one of three answers:
negative, zero, or positive. The \emph{zero}, or \emph{on-line},
answer means the hypothesis lies on the queried line, and its
answer class is $S_\ell^0=\{h\in S:f_\ell(h)=0\}$.
A \emph{ternary decision tree} has an outgoing edge for each
possible answer; a proper such tree retains only nonempty outcomes.
Leaf depths and both cost objectives use the same conventions as in
the binary model.
Use the two virtual binary predicates
\[
 b_1(h)=\mathbf1[f_\ell(h)<0],\qquad
 b_2(h)=\mathbf1[f_\ell(h)\le0].
\]
Their pair of bits distinguishes the three outcomes.
In the order $(b_1,b_2)$, negative, zero, and positive signs map to
$(1,1)$, $(0,1)$, and $(0,0)$, respectively; the unused pair $(1,0)$
never occurs. Thus no outcome is merged, and the virtual binary tree
can omit the unused branch at each query.
Requesting either bit is implemented by querying the original line
once and storing its full sign. Both bits are then known. For analysis
we count each requested virtual bit as a separate binary query;
the actual number of supplied-line queries can only be smaller.

For preprocessing, we also represent these bit labelings by planar
queries whose lines avoid every hypothesis. On the finite hypothesis
set, sufficiently small shifts preserve both bit values while moving
the virtual boundaries off the points. During execution, the algorithm
asks the original supplied line and reads its full ternary sign, so
the virtual representation changes only the analysis and preprocessing.
Any query whose answer is constant on the hypothesis set can be
discarded. For every remaining query, choose $\epsilon_\ell$ to be
half the minimum nonzero $|f_\ell(h)|$ on the finite hypothesis set.
At least one such nonzero value exists because the query is not
constant (a nonconstant sign labeling must contain a nonzero value).
For a rational input instance, this minimum is computable by scanning
the hypotheses and has polynomial encoding length: evaluating an
input affine form and taking the minimum of finitely many rational
absolute values preserves polynomial bit length. Then $b_1$ is represented
by $f_\ell(x)<-\epsilon_\ell$ and $b_2$ by
$f_\ell(x)<\epsilon_\ell$. These strict inequalities agree with the
definitions on every
hypothesis: a negative nonzero value is at most $-2\epsilon_\ell$,
and a positive value is at least $2\epsilon_\ell$, while zero lies
between the two thresholds. In particular, neither shifted boundary
contains a hypothesis. The shifts are used only for preprocessing;
their answers are obtained from the original supplied line.

\begin{symboltable}
$f_\ell(h)$ & Affine value whose sign is the answer for hypothesis $h$.\\
$b_1(h),\ b_2(h)$ & Virtual bits for $f_\ell(h)<0$ and $f_\ell(h)\le0$; together they encode the ternary sign.\\
$\epsilon_\ell$ & Positive shift smaller than every nonzero $|f_\ell(h)|$ on the finite hypothesis set.
\end{symboltable}

\begin{algorithm}[H]
\caption{Transfer a binary algorithm to on-line outcomes}
\label{alg:boundary}
\KwIn{Finite hypotheses, supplied ternary lines, binary tree-construction routine $B$.}
\KwOut{An executable tree using the original supplied lines.}
\ForEach{nonconstant supplied line $\ell$}{
  $\epsilon_\ell\gets\frac12\min_{h:f_\ell(h)\ne0}|f_\ell(h)|$\;
  Represent virtual bits $b_1=\mathbf1[f_\ell(h)<0]$ and
  $b_2=\mathbf1[f_\ell(h)\le0]$ by the two shifted off-boundary lines\;
}
Construct a virtual binary tree with $B$\;
Execute a requested virtual bit by querying its original line,
caching the entire ternary sign; answer further requests from that cache\;
\KwRet the contracted original-line tree\;
\end{algorithm}

Lines constant on the whole hypothesis set are discarded before taking
the minimum nonzero affine value. Proposition~\ref{prop:boundary}
proves that the approximation ratio increases by at most two.


There are two comparisons to make. An optimal ternary tree can be
simulated using at most two virtual bits per query. A constructed
virtual tree can then be executed using original line queries,
with the cached sign supplying both bits whenever either is asked.
Because the first comparison bounds each hypothesis's virtual depth
by twice its original depth, and the second bounds each executed depth
by that hypothesis's virtual depth, both comparisons survive either
taking a maximum or summing depths with nonnegative prior weights.
\begin{proposition}[Boundary reduction]\label{prop:boundary}
An $\alpha$-approximation for either binary objective gives a
$2\alpha$-approximation for the corresponding three-answer objective.
\end{proposition}
\begin{proof}
Simulate every query of an optimal original tree by its two virtual
bits. Each hypothesis uses at most twice as many binary queries,
so the virtual optimum is at most twice the original optimum
for either objective. Conversely, execute a constructed virtual
tree by original line queries and cached signs. Its actual depth
is no larger than its virtual depth on every hypothesis.
To combine these comparisons explicitly, write $C_{\rm orig}$ for
the executed cost and $\OPT_{\rm virt},\OPT_{\rm orig}$ for the two
optima. Then, for either depth objective,
\[
 C_{\rm orig}\le C_{\rm virt}
 \le\alpha\OPT_{\rm virt}\le2\alpha\OPT_{\rm orig}.
\]
The first step uses cached-sign execution, the second is the assumed
binary approximation guarantee, and the last uses the two-bit
simulation of the optimal ternary tree. Thus the comparison applies
to both maximum depth and any weighted sum of leaf depths.
\end{proof}

The reduction is independent of the binary algorithm's internal
geometry or schedule. Applying it to the strongest deterministic
algorithms established above gives the following bounds for the
model with on-line outcomes.
\begin{corollary}[Constant approximations with on-line outcomes]
\label{cor:constant-boundaries}
With on-line answers, expected depth has a polynomial-time
approximation factor below $10.365952$ for arbitrary positive priors,
and deterministic maximum depth has a factor below $8.929931$.
\end{corollary}
\begin{proof}
Apply Proposition~\ref{prop:boundary} to
Theorems~\ref{thm:two-phase-expected} and~\ref{thm:poly-slot-worst}.
\end{proof}
